\documentclass[11pt]{article}

\usepackage[french]{babel}

\usepackage[T1]{fontenc}

\usepackage{amsfonts}
\usepackage{amsmath}
\usepackage{amssymb}

\usepackage{graphicx}
\usepackage{multirow}
\usepackage{multicol}
\usepackage{xcolor}
\usepackage[shortlabels]{enumitem}

\usepackage[margin=2cm]{geometry}

\usepackage{mathtools}
\usepackage{siunitx}
\usepackage{physics}

\usepackage{tikz}
\usepackage{hyperref}

\usepackage{algorithm}
\usepackage{algpseudocode}
\usetikzlibrary{quantikz2,babel}

\tikzset{qcircuit/.style={
    column sep=.9cm,
    row sep=0.7cm,
    operator/.style={draw=black!80, line width=0.5pt, fill=gray!5, minimum height=0.7cm, minimum width=0.7cm},
    phase/.style={draw=black!80, fill=black!20, circle, minimum size=4pt},
    slice/.style={font=\footnotesize\itshape, color=gray!50, dashed}
}}


\DeclarePairedDelimiter\ceil{\lceil}{\rceil}
\DeclarePairedDelimiter\floor{\lfloor}{\rfloor}


\AtBeginDocument{\RenewCommandCopy\qty\SI}

\makeatletter
   \def\vhrulefill#1{\leavevmode\leaders\hrule\@height#1\hfill \kern\z@}
\makeatother

\newcommand{\moy}[1]{\langle #1 \rangle}
\newcommand{\pDer}[2]{\frac{\partial #1}{\partial #2}}
\newcommand{\pDerd}[2]{\frac{\partial^2 #1}{\partial #2^2}}

\newcounter{questioncounter}
\newenvironment{question}[1]{
\refstepcounter{questioncounter}
\noindent\rule{0.05\textwidth}{1pt}\,\,\raisebox{-.3\ht\strutbox}{\textbf{Question \thequestioncounter \hspace{5mm} #1}}\,\,\vhrulefill{1pt}\par\vspace{0.4cm}
}{\par}

%======================================================
%en-tete
%======================================================
\begin{document}
\selectlanguage{french}
\noindent{\Large Anthony Wroblewski} \hfill {\Large \today}\\
{\Large Nathan Brochu}

\par\vspace{5mm}
\begin{center}
   {\Large PHQ344 - Physique Statistique 1} \\ \vspace{5mm}
   {\LARGE\sffamily Devoir \#1} \\ \vspace{5mm}
\end{center}

%======================================================
%question 1
%======================================================
\begin{question}{Élimination du bruit}
On note $a = \text{signal total}$, \quad $a_s= \text{signal désiré}$, \quad $a_b = \text{bruit}$, 

On a que: 
$$\sqrt{(\Delta a_b)^2} \gg a_s \quad \text{où} \quad \Delta a_b= a_b + \langle a_b\rangle \quad et \quad  \langle a_b \rangle = 0 $$
\begin{enumerate}[label =(\alph*)]
   \item \textbf{Moyen du signal:}\\
Puisque le signal envoyer est le même à chaque mesure et que le bruit est aléatoire de moyen nulle. on peut répéter la mesure et faire la moyenne, pour faire diminuer le bruit.\\

On moyen le signal:
\begin{align*}
\langle a \rangle &= \frac{1}{N}\sum_i^N a_i\\
&=\frac{1}{N}\sum_i^N \left(a_s + a_b\right)\\
&= \frac{1}{N}\sum_i^N a_s + \frac{1}{N}\sum_i^N a_b\\
&=\langle a_s \rangle +\langle a_b \rangle
\end{align*}
On sais que $\langle a \rangle=0$ et puisque le signal $a_s$ est toujours le même, ca moyen est égale au signal. $\langle a_s \rangle=a_s$.\\
on a donc 
$$\langle a_s \rangle +\langle a_b \rangle = a_s + 0 = \boxed{a_s}$$
On récupère donc le signal original $a_s$

\item \textbf{Distinction du bruit}\\
on cherche a savoir de fois on doit envoyer un signal pour qu'il soit distingable du bruit. En d'autre mots, on cherche a avoir la relation entre on souhaite que l'amplitude du signal $a_s$ soit supérieure ou comparable
à l'incertitude associée au bruit moyen après $N$ mesures :\\

Pour pouvoir distinguer le signal du bruit, on souhaite avoir

\[
\boxed{
a_s \gtrsim \Delta\overline{a}_b
}.
\]

Par définition,

\begin{align*}
\Delta\overline{a}_b &= \sqrt{\operatorname{Var}(\overline{a}_b)}.\\
\operatorname{Var}(\overline{a}_b) &= \operatorname{Var}\left(\frac{1}{N}\sum_{i=1}^{N}a_{b,i}\right).\\
\operatorname{Var}(\overline{a}_b) &= \frac{1}{N^2}\operatorname{Var}\left(\sum_{i=1}^{N}a_{b,i}\right).\\
\operatorname{Var}(\overline{a}_b) &= \frac{1}{N^2}\sum_{i=1}^{N}\operatorname{Var}(a_{b,i}).
\end{align*}

Puisque

\[
\operatorname{Var}(a_{b,i})
=
(\Delta a_b)^2,
\]

on obtient

\[
\operatorname{Var}(\overline{a}_b)
=
\frac{1}{N^2}
N(\Delta a_b)^2
=
\frac{(\Delta a_b)^2}{N}.
\]

Ainsi,

\[
\boxed{
\Delta\overline{a}_b
=
\frac{\Delta a_b}{\sqrt{N}}
}.
\]

On remplace ce résultat dans la condition permettant de distinguer
le signal du bruit :

\[
a_s
\gtrsim
\frac{\Delta a_b}{\sqrt{N}}.
\]

On isole alors $N$ :

\[
a_s\sqrt{N}
\gtrsim
\Delta a_b,
\]

\[
\sqrt{N}
\gtrsim
\frac{\Delta a_b}{a_s},
\]

et finalement,

\[
\boxed{
N
\gtrsim
\frac{(\Delta a_b)^2}{a_s^2}
}.
\]

Ainsi, le nombre de mesures nécessaires correspond approximativement
au rapport entre l'intensité du bruit $(\Delta a_b)^2$ et l'intensité
du signal $a_s^2$.
\end{enumerate}
\end{question}

%======================================================
%question 2
%======================================================
\begin{question}{Configurations de bosons}
\begin{enumerate}[label=(\alph*)]
   \item 

On considère $n$ barres verticales identiques et $p$ objets
identiques. On possède donc au total

\[
n+p
\]

éléments à ordonner.

Si tous les éléments étaient différents, il y aurait

\[
(n+p)!
\]

façons de les ordonner. Cependant, les $p$ objets sont identiques, donc
les $p!$ permutations entre les objets ne donnent pas de nouveaux
arrangements. De même, les $n$ barres identiques. 

Le nombre d'arrangements distincts est donc

\[
\boxed{
\Omega
=
\frac{(n+p)!}{n!p!}
}.
\]

On peut également écrire

\[
\boxed{
\Omega =
\binom{n+p}{p} = \frac{(n+p)!}{p!((n+p)-p)!} = \frac{(n+p)!}{p!n!}
=
\binom{n+p}{n} 
%= \frac{(n+p)!}{n!((p+n)-n)!}
}.
\]

\bigskip

\item On considère maintenant $N$ particules identiques que l'on
souhaite répartir dans $m$ boîtes. Chaque boite peut contenir un nombre
arbitraire de particules.

On représente les $N$ particules par des objets identiques. Pour séparer
les $m$ boîtes, on a besoin de

\[
m-1
\]

barres verticales.\\
exemple pour 5 boite, 10 objets : $\underbrace{ooo|}_{boite 1}oo\underbrace{|\hspace{0.2 cm}|}_{boite 3}ooo\underbrace{|oo}_{boite 5}$\\

Le problème revient donc à ordonner $N$ objets identiques et $m-1$
barres identiques.

En utilisant le résultat de la partie A, on pose

\[
p=N
\qquad \text{et} \qquad
n=m-1.
\]

On obtient donc

\[
\Omega
=
\frac{(N+m-1)!}{N!(m-1)!}.
\]

Le nombre de façons de répartir $N$ bosons identiques dans $m$
niveaux d'énergie est

\[
\boxed{
\Omega_B
=
\frac{(N+m-1)!}{N!(m-1)!}
=
\binom{N+m-1}{N}
}.
\]

% \textbf{Intuition :} les bosons sont identiques et plusieurs bosons
% peuvent occuper le même niveau d'énergie. Les $m-1$ barres servent donc
% simplement à séparer les $N$ particules entre les $m$ niveaux.
\end{enumerate}

\end{question}

%======================================================
%question 3
%======================================================
\begin{question}{Configurations de fermions}

On considère $N$ fermions identiques que l'on souhaite répartir dans
$M$ niveaux d'énergie, avec

\[
M>N.
\]

Contrairement aux bosons, un niveau d'énergie ne peut contenir qu'un
seul fermion. Chaque niveau possède donc seulement deux possibilités :

\[
\text{vide}
\qquad \text{ou} \qquad
\text{occupé}.
\]

Comme les fermions sont identiques, leur ordre n'a pas d'importance.
Il suffit donc de choisir quels $N$ niveaux parmi les $M$ niveaux
disponibles seront occupés.

Le nombre de façons de choisir $N$ niveaux parmi $M$ est donné par

\[
\binom{M}{N}.
\]

On obtient donc

\[
\boxed{
\Omega_F
=
\binom{M}{N}
=
\frac{M!}{N!(M-N)!}
}.
\]

% \textbf{Intuition :} contrairement aux bosons, plusieurs fermions ne
% peuvent pas occuper le même niveau d'énergie. Le problème revient donc
% simplement à choisir $N$ niveaux occupés parmi les $M$ niveaux
% disponibles.
\end{question}



\end{document}